\subsection{应用：群的上同调与同调}

设 G 是一个群， \(\mathbf{Z}^{(G)}\) 为其上的 Z-代数（III，第 19 页）。回顾（参见 III，第 20 页，例）若 M 是一个交换群，则给出 G 在 M 上的作用（即一个同态 \(\tau: G \to \text{Aut}(M)\) ），与给出加法群 M 上的 \(\mathbf{Z}^{(G)}\) -左模结构等价。特别地，我们将把群 Z 视为一个 \(\mathbf{Z}^{(G)}\) -左模，并赋予其平凡作用。

定义 1. --- 设 M 为一个 \(\mathbf{Z}^{(G)}\) -左（相应地，右）模， \(n\) 为一个整数 \(\geqslant 0\) 。群 \(\operatorname{Ext}_{\mathbf{Z}^{(G)}}^{n}(\mathbf{Z}, \mathbf{M})\) (相应地， \(\operatorname{Tor}_{n}^{\mathbf{Z}^{(G)}}(\mathbf{M}, \mathbf{Z})\) ) 记为 \(\mathrm{H}^{n}(\mathrm{G}, \mathbf{M})\) (相应地， \(\mathrm{H}_{n}(\mathrm{G}, \mathbf{M}))\) 并称为 \(n\) 次上同调群（相应地，同调群），即 \(\mathbf{G}\) 的以 \(\mathbf{M}\) 为系数的上同调群（相应地，同调群）。

标准分解（X，第 58 页）B(Z \(^{(G)}\) ，Z) 是 Z \(^{(G)}\) -模 Z 的一个自由分解；由此可知，群 H \(^{n}\) (G, M)（相应地 H \(_{n}\) (G, M)）与以下复形的同调群等同：

\[
\operatorname{Hom} _ {\mathbf {Z} ^ {(G)}} \left(\mathrm{B} (\mathbf {Z} ^ {(G)}, \mathbf {Z}), \mathrm{M}\right) \quad \left(\text { resp. } \mathrm{M} \otimes_ {\mathbf {Z} ^ {(G)}} \mathrm{B} (\mathbf {Z} ^ {(G)}, \mathbf {Z})\right).
\]

利用 \((\mathbf{Z}^{(\mathbf{G})})^{\otimes n}\) 到 \(\mathbf{Z}^{(\mathbf{G}^{n})}\) （III，第 36 页）上的典范同构以及纯量扩张的性质（II，第 82 页），我们得出结论： \(\mathrm{H}^n (\mathrm{G},\mathrm{M})\) 典范同构于 \(n\) 次上升同调群，相应的复形 \(\mathrm{C(G,M)}\) 按以下方式定义： \(\mathbf{C}^n (\mathbf{G},\mathbf{M}) = 0\) 当 \(n < 0\) ；设 \(n\geqslant 0\) ， \(\mathbf{C}^n (\mathbf{G},\mathbf{M})\) 是 \(\mathbf{Z}\) -模，其元素为从 \(\mathbf{G}^n\) 到 \(\mathbf{M}\) 的映射；设 \(n\geqslant 0\) ，微分 \(d^n:\mathbf{C}^n (\mathbf{G},\mathbf{M})\to \mathbf{C}^{n + 1}(\mathbf{G},\mathbf{M})\) 由以下给出

\[
(d ^ {n} f) (g _ {0}, \dots , g _ {n}) = g _ {0} \cdot f (g _ {1}, \dots , g _ {n}) + \sum_ {i = 0} ^ {n - 1} (- 1) ^ {i + 1} f (g _ {0}, \dots , g _ {i} g _ {i + 1}, \dots , g _ {n})
\]

\[
+ (- 1) ^ {n + 1} f \left(g _ {0}, \dots , g _ {n - 1}\right)
\]

其中 \(f\) 属于 \(\mathbf{C}^{n}(\mathbf{G},\mathbf{M})\) ， \(g_{0},\ldots,g_{n}\) 属于 \(\mathbf{G}\) 。

同样地， \(\mathrm{H}_{n}(\mathrm{G},\mathrm{M})\) 与复形 \(\mathrm{C}'(\mathrm{G},\mathrm{M})\) 的 n 阶同调群等同，其中 \(\mathrm{C}_{n}^{\prime}(\mathrm{G},\mathrm{M})=\mathrm{M}\otimes_{\mathbf{Z}}\mathbf{Z}^{(\mathbf{G}^{n})}\) 当 \(n\geqslant0\) ， \(\mathrm{C}_{n}^{\prime}(\mathrm{G},\mathrm{M})=0\) 当 n\textless0 ，微分 \(d_{n}:C_{n}^{\prime}(\mathrm{G},\mathrm{M})\to C_{n-1}^{\prime}(\mathrm{G},\mathrm{M})\) 由以下定义：

\[
d _ {n} (m \otimes e _ {g _ {1}, \dots , g _ {n}}) = m. g _ {1} \otimes e _ {g _ {2}, \dots , g _ {n}} + \sum_ {i = 1} ^ {n - 1} (- 1) ^ {i} m \otimes e _ {g _ {1}, \dots , g _ {i} g _ {i + 1}, \dots , g _ {n}}
\]

\[
+ (- 1) ^ {n} m \otimes e _ {g _ {1}, \dots , g _ {n - 1}}
\]

其中 \(n \geqslant 1\) , \(m\) 属于 M，且 \(g_1, \ldots, g_n\) 属于 G。

例。--- 1) 由定义直接可得 H \(^{0}\) (G, M) 同构于 M 中在 G 作用下不变的元素所成的子模，且 H \(_{0}\) (G, M) 同构于 M 的商模，即 M 除以由元素 m.g - m（其中 m ∈ M，g ∈ G）生成的子模所得的商。

\begin{enumerate}
\def\labelenumi{\arabic{enumi})}
\setcounter{enumi}{1}
\tightlist
\item
  由此可知 \(H^{1}(G, M)\) 同构于 Z-模 \(Z^{1}(G, M)/B^{1}(G, M)\) ，其中 \(Z^{1}(G, M)\) 是满足以下条件的从 G 到 M 的函数 f 所构成的 Z-模：
\end{enumerate}

\[
f \left(g _ {1}, g _ {2}\right) = g _ {1} \cdot f \left(g _ {2}\right) + f \left(g _ {1}\right) \quad \text { for   all } \quad g _ {1}, g _ {2} \text { in } G,
\]

且 \(\mathbf{B}^{1}(\mathbf{G},\mathbf{M})\) 是一个子 \(\mathbf{Z}\) -模，包含于 \(\mathbf{Z}^{1}(\mathbf{G},\mathbf{M})\) ，由满足以下条件的 \(f\) 组成：存在一个元素 \(m\) 属于 \(\mathbf{M}\) ，使得：

\[
f (g) = g. m - m \quad \text { for   every } \quad g \in G.
\]

有时人们说 \(Z^{1}(G, M)\) 是从 G 到 M 的交叉同态的 Z-模，且 \(B^{1}(G, M)\) 是主交叉同态的子模。

我们用 \(\tau: \mathbf{G} \to \operatorname{Aut}(\mathbf{M})\) 表示由 \(\mathbf{G}\) 的作用推出的同态；考虑外部半直积 \(\mathbf{M} \times_{\tau} \mathbf{G}\) 以及扩张 \(\xi_{\tau}: \mathbf{M} \xrightarrow{i} \mathbf{M} \times_{\tau} \mathbf{G} \xrightarrow{p} \mathbf{G}\) （I，第 64 页）。设 \(e: \mathbf{G} \to \mathbf{M} \times_{\tau} \mathbf{G}\) 为一个映射使得 \(p \circ e = 1_{\mathbf{G}}\) ；我们有 \(e = (f, 1_{\mathbf{G}})\) ，其中 \(f \in C^{1}(\mathbf{G}, \mathbf{M})\) 。 \(e\) 为同态（即扩张 \(\xi_{\tau}\) 的一个截面）当且仅当 \(f \in Z^{1}(\mathbf{G}, \mathbf{M})\) 。 \(\xi_{\tau}\) 的两个截面被 \(i(\mathbf{M})\) 的一个元素共轭，当且仅当相应的交叉同态在 \(H^{1}(\mathbf{G}, \mathbf{M})\) 中具有相同的类。

当 G 平凡地作用于 M 时，我们有 \(\mathbf{B}^{1}(\mathbf{G},\mathbf{M})=0\) 且 \(\mathbf{H}^{1}(\mathbf{G},\mathbf{M})\) 同构于从 G 到 M 的群同态所构成的 Z-模。

\begin{enumerate}
\def\labelenumi{\arabic{enumi})}
\setcounter{enumi}{2}
\tightlist
\item
  类似地， \(H^{2}(G, M)\) 同构于 Z-模 \(Z^{2}(G, M)/B^{2}(G, M)\) ，其中 \(Z^{2}(G, M)\) 是从 G × G 到 M 的映射 f 所构成的 Z-模，满足：
\end{enumerate}

\[
g _ {1} \cdot f (g _ {2}, g _ {3}) - f (g _ {1} g _ {2}, g _ {3}) + f (g _ {1}, g _ {2} g _ {3}) - f (g _ {1}, g _ {2}) = 0
\]

其中 \(g_{1}, g_{2}, g_{3}\) 属于 \(\mathbf{G}\) ，且 \(\mathbf{B}^{2}(\mathbf{G}, \mathbf{M})\) 是一个子 \(\mathbf{Z}\) -模，包含于 \(\mathbf{Z}^{2}(\mathbf{G}, \mathbf{M})\) ，由满足以下条件的 \(f\) 组成：存在映射 \(h\) 自 \(\mathbf{G}\) 到 \(\mathbf{M}\) ，使得：

\[
f (g _ {1}, g _ {2}) = g _ {1}. h (g _ {2}) - h (g _ {1} g _ {2}) + h (g _ {1})
\]

其中 \(g_{1}, g_{2}\) 属于 G。

因此我们重新得到了第八篇附录中给出的群 \(H^{2}(G, M)\) 的定义。特别地，它推出存在 \(H^{2}(G, M)\) 到 G 关于 M 的扩张类群上的一个典范同构（同上引）。

\begin{enumerate}
\def\labelenumi{\arabic{enumi})}
\setcounter{enumi}{3}
\tightlist
\item
  设 M 是一个 Z-模，我们将其视为一个 \(\mathbf{Z}^{(\mathrm{G})}\) -右模，令 G 平凡作用。群 \(\mathrm{H}_{1}(\mathrm{G},\mathrm{M})\) 同构于 \(\mathrm{M}\otimes_{\mathbf{Z}}\mathbf{Z}^{(\mathrm{G})}\) 的商，其中商为除以由元素 \(m\otimes(e_{g_{1}g_{2}}-e_{g_{1}}-e_{g_{2}})\) （m 属于 M， \(g_{1}, g_{2}\) 属于 G）生成的子 Z-模；由此容易推出 \(\mathrm{H}_{1}(\mathrm{G},\mathrm{M})\) 同构于 \(\mathrm{M}\otimes_{\mathbf{Z}}(\mathrm{G}/(\mathrm{G},\mathrm{G}))\) .
\end{enumerate}

我们用 \(\sigma\) 表示 \(\mathbf{Z}^{(G)}\) 的反自同构，它由 \(\sigma(e_g) = e_{g-1}\) 对 \(g \in G\) 定义。每个左 \(\mathbf{Z}^{(G)}\) -模都可以视为一个 \(\mathbf{Z}^{(G)}\) -右模（借助 \(\sigma\) ），反之亦然。这使我们能够，例如，定义群 \(\mathrm{H}_q(\mathrm{G}, \mathrm{M})\) （对左 \(\mathbf{Z}^{(G)}\) -模 M），通过设定 \(\mathrm{H}_q(\mathrm{G}, \mathrm{M}) = \mathrm{H}_q(\mathrm{G}, \sigma_*(\mathrm{M})) = \mathrm{Tor}_q^{\mathbf{Z}^{(G)}}(\mathbf{Z}, \mathbf{M})\) .

引理 1. --- 设 M 是一个 Z-模；记 M \(^{G}\) 为群 Hom \(_{Z}\) (Z \(^{(G)}\) , M) 并赋予其自然的左 Z \(^{(G)}\) -模结构。则：

\[
\mathrm{H} ^ {i} (\mathrm{G}, \mathrm{M} ^ {\mathrm{G}}) = 0 \quad \text { for } \quad i \geqslant 1.
\]

事实上，由命题 8，b)（X，第 109 页），取 A = N = Z \(^{(G)}\) 且 B = Q = Z，可得典范同构：

\[
\operatorname{Ext} _ {\mathbf {Z} ^ {(G)}} \left(\mathbf {Z}, \mathrm{M} ^ {\mathrm{G}}\right)\rightarrow \operatorname{Ext} _ {\mathbf {Z}} (\mathbf {Z}, \mathrm{M})
\]

由此即得引理。

命题 11. --- 设 L 是交换域 K 的有限伽罗瓦扩张，伽罗瓦群为 G。

\begin{enumerate}
\def\labelenumi{\alph{enumi})}
\item
  我们有 \(\mathrm{H}^i (\mathbf{G},\mathbf{L}) = 0\) 当 \(i\geqslant 1\)
\item
  我们有 \(\mathrm{H}^1 (\mathbf{G},\mathbf{L}^*) = 0\)
\item
  群 \(\mathrm{H}^2 (\mathbf{G},\mathbf{L}^*)\) 典范同构于群 Br(K, L)（VIII，§ 13）。
\end{enumerate}

正规基定理（V，§10，第 9 号，定理 6）表明 L 作为 \(\mathbf{Z}^{(G)}\) -模同构于 \(\mathbf{K}^{\mathrm{G}} = \operatorname{Hom}_{\mathbf{Z}}(\mathbf{Z}^{(G)}, \mathbf{K})\) ；断言 \(a)\) 随后由引理 1 得出。鉴于例 2，断言 \(b)\) 由 V，§10，第 5 号，命题 9 的推论 1 得出；最后，断言 \(c)\) 已在 VIII，§13 中证明。

